\section{Lemmas on bivariant classes}
\label{section-bivariant-II}

\noindent
In this section we prove some elementary results on bivariant classes.
Here is a criterion to see that an operation
passes through rational equivalence.

\begin{lemma}
\label{lemma-factors-through-rational-equivalence}
\begin{reference}
Very weak form of \cite[Theorem 17.1]{F}
\end{reference}
Let $(S, \delta)$ be as in Situation \ref{situation-setup}.
Let $f : X \to Y$ be a morphism of schemes locally of finite type over $S$.
Let $p \in \mathbf{Z}$. Suppose given a rule
which assigns to every locally of finite type morphism $Y' \to Y$
and every $k$ a map
$$
c \cap - : Z_k(Y') \longrightarrow \CH_{k - p}(X')
$$
where $Y' = X' \times_X Y$, satisfying condition (3) of
Definition \ref{definition-bivariant-class}
whenever $\mathcal{L}'|_{D'} \cong \mathcal{O}_{D'}$. Then
$c \cap -$ factors through rational equivalence.
\end{lemma}

\begin{proof}
The statement makes sense because given a triple
$(\mathcal{L}, s, i : D \to X)$ as in
Definition \ref{definition-gysin-homomorphism}
such that $\mathcal{L}|_D \cong \mathcal{O}_D$, then
the operation $i^*$ is defined on the level of cycles, see
Remark \ref{remark-gysin-on-cycles}.
Let $\alpha \in Z_k(X')$ be a cycle which is rationally equivalent to zero.
We have to show that $c \cap \alpha = 0$. By
Lemma \ref{lemma-rational-equivalence-family}
there exists a cycle $\beta \in Z_{k + 1}(X' \times \mathbf{P}^1)$
such that $\alpha = i_0^*\beta - i_\infty^*\beta$
where $i_0, i_\infty : X' \to X' \times \mathbf{P}^1$ are the
closed immersions of $X'$ over $0, \infty$. Since these are
examples of effective Cartier divisors with trivial normal
bundles, we see that $c \cap i_0^*\beta = j_0^*(c \cap \beta)$
and $c \cap i_\infty^*\beta = j_\infty^*(c \cap \beta)$
where $j_0, j_\infty : Y' \to Y' \times \mathbf{P}^1$ are
closed immersions as before. Since
$j_0^*(c \cap \beta) \sim_{rat} j_\infty^*(c \cap \beta)$
(follows from Lemma \ref{lemma-rational-equivalence-family}) we conclude.
\end{proof}

\begin{lemma}
\label{lemma-bivariant-weaker}
\begin{reference}
Weak form of \cite[Theorem 17.1]{F}
\end{reference}
Let $(S, \delta)$ be as in Situation \ref{situation-setup}.
Let $f : X \to Y$ be a morphism of schemes locally of finite type over $S$.
Let $p \in \mathbf{Z}$. Suppose given a rule
which assigns to every locally of finite type morphism $Y' \to Y$
and every $k$ a map
$$
c \cap - : \CH_k(Y') \longrightarrow \CH_{k - p}(X')
$$
where $Y' = X' \times_X Y$, satisfying conditions (1), (2) of
Definition \ref{definition-bivariant-class}
and condition (3) whenever $\mathcal{L}'|_{D'} \cong \mathcal{O}_{D'}$. Then
$c \cap -$ is a bivariant class.
\end{lemma}

\begin{proof}
Let $Y' \to Y$ be a morphism of schemes which is locally of finite type.
Let $(\mathcal{L}', s', i' : D' \to Y')$ be as in
Definition \ref{definition-gysin-homomorphism}
with pullback $(\mathcal{N}', t', j' : E' \to X')$ to $X'$.
We have to show that $c \cap (i')^*\alpha' = (j')^*(c \cap \alpha')$
for all $\alpha' \in \CH_k(Y')$.

\medskip\noindent
Denote $g : Y'' \to Y'$ the smooth morphism of relative
dimension $1$ with $i'' : D'' \to Y''$ and $p : D'' \to D'$
constructed in Lemma \ref{lemma-normal-cone-effective-Cartier}.
(Warning: $D''$ isn't the full inverse image of $D'$.)
Denote $f : X'' \to X'$ and $E'' \subset X''$
their base changes by $X' \to Y'$. Picture
$$
\xymatrix{
& X'' \ar[rr] \ar'[d][dd]_h & & Y'' \ar[dd]^g \\
E'' \ar[rr] \ar[dd]_q \ar[ru]^{j''} & & D'' \ar[dd]^p \ar[ru]^{i''} & \\
& X' \ar'[r][rr] & & Y' \\
E' \ar[rr] \ar[ru]^{j'} & & D' \ar[ru]^{i'}
}
$$
By the properties given in the lemma we know that $\beta' = (i')^*\alpha'$
is the unique element of $\CH_{k - 1}(D')$ such that
$p^*\beta' = (i'')^*g^*\alpha'$. Similarly, we know that
$\gamma' = (j')^*(c \cap \alpha')$ is the unique element of
$\CH_{k - 1 - p}(E')$ such that $q^*\gamma' = (j'')^*h^*(c \cap \alpha')$.
Since we know that
$$
(j'')^*h^*(c \cap \alpha') =
(j'')^*(c \cap g^*\alpha') =
c \cap (i'')^*g^*\alpha'
$$
by our assumptions on $c$; note that the modified version of (3)
assumed in the statement of the lemma applies to $i''$
and its base change $j''$. We similarly know that
$$
q^*(c \cap \beta') = c \cap p^*\beta'
$$
We conclude that $\gamma' = c \cap \beta'$ by the uniqueness pointed
out above.
\end{proof}

\noindent
Here a criterion for when a bivariant class is zero.

\begin{lemma}
\label{lemma-bivariant-zero}
Let $(S, \delta)$ be as in Situation \ref{situation-setup}.
Let $f : X \to Y$ be a morphism of schemes locally of finite type over $S$.
Let $c \in A^p(X \to Y)$. For $Y'' \to Y' \to Y$ set
$X'' = Y'' \times_Y X$ and $X' = Y' \times_Y X$.
The following are equivalent
\begin{enumerate}
\item $c$ is zero,
\item $c \cap [Y'] = 0$ in $\CH_*(X')$ for every integral scheme $Y'$
locally of finite type over $Y$, and
\item for every integral scheme $Y'$ locally of finite type over $Y$,
there exists a proper birational morphism $Y'' \to Y'$ such that
$c \cap [Y''] = 0$ in $\CH_*(X'')$.
\end{enumerate}
\end{lemma}

\begin{proof}
The implications (1) $\Rightarrow$ (2) $\Rightarrow$ (3) are clear.
Assumption (3) implies (2) because $(Y'' \to Y')_*[Y''] = [Y']$
and hence $c \cap [Y'] = (X'' \to X')_*(c \cap [Y''])$ as $c$
is a bivariant class. Assume (2).
Let $Y' \to Y$ be locally of finite type. Let $\alpha \in \CH_k(Y')$.
Write $\alpha = \sum n_i [Y'_i]$ with $Y'_i \subset Y'$ a locally finite
collection of integral closed subschemes of $\delta$-dimension $k$.
Then we see that $\alpha$ is pushforward of the cycle
$\alpha' = \sum n_i[Y'_i]$ on $Y'' = \coprod Y'_i$ under the
proper morphism $Y'' \to Y'$. By the properties of bivariant
classes it suffices to prove that $c \cap \alpha' = 0$ in $\CH_{k - p}(X'')$.
We have $\CH_{k - p}(X'') = \prod \CH_{k - p}(X'_i)$ where
$X'_i = Y'_i \times_Y X$. This follows immediately
from the definitions. The projection maps
$\CH_{k - p}(X'') \to \CH_{k - p}(X'_i)$ are given by flat pullback.
Since capping with $c$ commutes with
flat pullback, we see that it suffices to show that $c \cap [Y'_i]$
is zero in $\CH_{k - p}(X'_i)$ which is true by assumption.
\end{proof}

\begin{lemma}
\label{lemma-disjoint-decomposition-bivariant}
Let $(S, \delta)$ be as in Situation \ref{situation-setup}.
Let $f : X \to Y$ be a morphism of schemes locally of finite type over $S$.
Assume we have disjoint union decompositions
$X = \coprod_{i \in I} X_i$ and $Y = \coprod_{j \in J} Y_j$
by open and closed subschemes
and a map $a : I \to J$ of sets such that $f(X_i) \subset Y_{a(i)}$.
Then
$$
A^p(X \to Y) = \prod\nolimits_{i \in I} A^p(X_i \to Y_{a(i)})
$$
\end{lemma}

\begin{proof}
Suppose given an element $(c_i) \in \prod_i A^p(X_i \to Y_{a(i)})$.
Then given $\beta \in \CH_k(Y)$ we can map this to the element of
$\CH_{k - p}(X)$ whose restriction to $X_i$ is $c_i \cap \beta|_{Y_{a(i)}}$.
This works because $\CH_{k - p}(X) = \prod_i \CH_{k - p}(X_i)$.
The same construction works after base change by any $Y' \to Y$
locally of finite type and we get $c \in A^p(X \to Y)$.
Thus we obtain a map $\Psi$ from the right hand side of the formula
to the left hand side of the formula.
Conversely, given $c \in A^p(X \to Y)$ and an element
$\beta_i \in \CH_k(Y_{a(i)})$ we can consider the element
$(c \cap (Y_{a(i)} \to Y)_*\beta_i)|_{X_i}$ in $\CH_{k - p}(X_i)$.
The same thing works after base change by any $Y' \to Y$
locally of finite type and we get $c_i \in A^p(X_i \to Y_{a(i)})$.
Thus we obtain a map $\Phi$ from the left hand
side of the formula to the right hand side of the formula.
It is immediate that $\Phi \circ \Psi = \text{id}$.
For the converse, suppose that $c \in A^p(X \to Y)$ and
$\beta \in \CH_k(Y)$. Say $\Phi(c) = (c_i)$. Let $j \in J$.
Because $c$ commutes with flat pullback we get
$$
(c \cap \beta)|_{\coprod_{a(i) = j} X_i} =
c \cap \beta|_{Y_j}
$$
Because $c$ commutes with proper pushforward we get
$$
(\coprod\nolimits_{a(i) = j} X_i \to X)_*
((c \cap \beta)|_{\coprod_{a(i) = j} X_i})
=
c \cap (Y_j \to Y)_*\beta|_{Y_j}
$$
The left hand side is the cycle on $X$ restricting to $(c \cap \beta)|_{X_i}$
on $X_i$ for $i \in I$ with $a(i) = j$ and $0$ else.
The right hand side is a cycle on $X$ whose restriction to $X_i$
is $c_i \cap \beta|_{Y_j}$ for $i \in I$ with $a(i) = j$.
Thus $c \cap \beta = \Psi((c_i))$ as desired.
\end{proof}

\begin{remark}
\label{remark-completion-bivariant}
Let $(S, \delta)$ be as in Situation \ref{situation-setup}.
Let $f : X \to Y$ be a morphism of schemes locally of finite type over $S$.
Let $X = \coprod_{i \in I} X_i$ and $Y = \coprod_{j \in J} Y_j$
be the decomposition of $X$ and $Y$ into their connected components
(the connected components are open as $X$ and $Y$ are locally Noetherian, see
Topology, Lemma \ref{topology-lemma-locally-Noetherian-locally-connected} and
Properties, Lemma \ref{properties-lemma-Noetherian-topology}).
Let $a(i) \in J$ be the index such that $f(X_i) \subset Y_{a(i)}$.
Then $A^p(X \to Y) = \prod A^p(X_i \to Y_{a(i)})$ by
Lemma \ref{lemma-disjoint-decomposition-bivariant}.
In this setting it is convenient to set
$$
A^*(X \to Y)^\wedge = \prod\nolimits_i A^*(X_i \to Y_{a(i)})
$$
This ``completed'' bivariant group is the subset
$$
A^*(X \to Y)^\wedge \quad\subset\quad \prod\nolimits_{p \geq 0} A^p(X \to Y)
$$
consisting of elements $c = (c_0, c_1, c_2, \ldots)$ such that
for each connected component $X_i$ the image of $c_p$ in
$A^p(X_i \to Y_{a(i)})$ is zero for almost all $p$.
If $Y \to Z$ is a second morphism, then the
composition $A^*(X \to Y) \times A^*(Y \to Z) \to A^*(X \to Z)$
extends to a composition
$A^*(X \to Y)^\wedge \times A^*(Y \to Z)^\wedge \to A^*(X \to Z)^\wedge$
of completions. 
We sometimes call $A^*(X)^\wedge = A^*(X \to X)^\wedge$ the
{\it completed bivariant cohomology ring} of $X$.
\end{remark}

\begin{lemma}
\label{lemma-envelope-bivariant}
Let $(S, \delta)$ be as in Situation \ref{situation-setup}.
Let $f : X \to Y$ be a morphism of schemes locally of finite type over $S$.
Let $g : Y' \to Y$ be an envelope (Definition \ref{definition-envelope})
and denote $X' = Y' \times_Y X$. Let $p \in \mathbf{Z}$ and let
$c' \in A^p(X' \to Y')$. If the two restrictions
$$
res_1(c') = res_2(c') \in A^p(X' \times_X X' \to Y' \times_Y Y')
$$
are equal (see proof), then there exists a unique $c \in A^p(X \to Y)$
whose restriction $res(c) = c'$ in $A^p(X' \to Y')$.
\end{lemma}

\begin{proof}
We have a commutative diagram
$$
\xymatrix{
X' \times_X X' \ar[d]^{f''} \ar@<1ex>[r]^-a \ar@<-1ex>[r]_-b &
X' \ar[d]^{f'} \ar[r]_h &
X \ar[d]^f \\
Y' \times_Y Y' \ar@<1ex>[r]^-p \ar@<-1ex>[r]_-q &
Y' \ar[r]^g &
Y
}
$$
The element $res_1(c')$ is the restriction (see
Remark \ref{remark-restriction-bivariant}) of $c'$ for the cartesian square
with morphisms $a, f', p, f''$ and the element $res_2(c')$ is the restriction
of $c'$ for the cartesian square with morphisms $b, f', q, f''$.
Assume $res_1(c') = res_2(c')$ and let $\beta \in \CH_k(Y)$.
By Lemma \ref{lemma-envelope} we can find a $\beta' \in \CH_k(Y')$
with $g_*\beta' = \beta$. Then we set
$$
c \cap \beta = h_*(c' \cap \beta')
$$
To see that this is independent of the choice of
$\beta'$ it suffices to show that
$h_*(c' \cap (p_*\gamma - q_*\gamma))$ is zero
for $\gamma \in \CH_k(Y' \times_Y Y')$.
Since $c'$ is a bivariant class we have
$$
h_*(c' \cap (p_*\gamma - q_*\gamma)) =
h_*(a_*(c' \cap \gamma) - b_*(c' \cap \gamma)) = 0
$$
the last equality since $h_* \circ a_* = h_* \circ b_*$
as $h \circ a = h \circ b$.

\medskip\noindent
Observe that our choice for $c \cap \beta$ is forced
by the requirement that $res(c) = c'$ and the compatibility
of bivariant classes with proper pushforward.

\medskip\noindent
Of course, in order to define the bivariant class $c$ we need
to construct maps $c \cap -: \CH_k(Y_1) \to \CH_{k + p}(Y_1 \times_Y X)$
for any morphism $Y_1 \to Y$ locally of finite type satisfying the
conditions listed in Definition \ref{definition-bivariant-class}.
Denote $Y'_1 = Y' \times_Y Y_1$, $X_1 = X \times_Y Y_1$.
The morphism $Y'_1 \to Y_1$ is an envelope by
Lemma \ref{lemma-base-change-envelope}. Hence we can use the base
changed diagram
$$
\xymatrix{
X'_1 \times_{X_1} X'_1 \ar[d]^{f''_1} \ar@<1ex>[r]^-{a_1} \ar@<-1ex>[r]_-{b_1} &
X'_1 \ar[d]^{f'_1} \ar[r]_{h_1} &
X_1 \ar[d]^{f_1} \\
Y'_1 \times_{Y_1} Y'_1 \ar@<1ex>[r]^-{p_1} \ar@<-1ex>[r]_-{q_1} &
Y'_1 \ar[r]^{g_1} &
Y_1
}
$$
and the same arguments to get a well defined map
$c \cap - : \CH_k(Y_1) \to \CH_{k + p}(X_1)$ as before.

\medskip\noindent
Next, we have to check conditions (1), (2), and (3) of
Definition \ref{definition-bivariant-class} for $c$.
For example, suppose that $t : Y_2 \to Y_1$ is a proper morphism
of schemes locally of finite type over $Y$. Denote as above
the base changes of the first diagram to $Y_1$, resp.\ $Y_2$,
by subscripts ${}_1$, resp.\ ${}_2$. Denote $t' : Y'_2 \to Y'_1$,
$s : X_2 \to X_1$, and $s' : X'_2 \to X'_1$ the base changes of $t$
to $Y'$, $X$, and $X'$. We have to show that
$$
s_*(c \cap \beta_2) = c \cap t_*\beta_2
$$
for $\beta_2 \in \CH_k(Y_2)$. Choose $\beta'_2 \in \CH_k(Y'_2)$
with $g_{2, *}\beta'_2 = \beta_2$. Since $c'$ is a bivariant class
and the diagrams
$$
\vcenter{
\xymatrix{
X'_2 \ar[d]_{s'} \ar[r]_{h_2} & X_2 \ar[d]^s \\
X'_1 \ar[r]^{h_1} & X_1
}
}
\quad\text{and}\quad
\vcenter{
\xymatrix{
X'_2 \ar[d]_{s'} \ar[r]_{f'_2} & Y'_2 \ar[d]^{t'} \\
X'_2 \ar[r]^{f'_1} & Y'_1
}
}
$$
are cartesian we have
$$
s_*(c \cap \beta_2) =
s_*(h_{2, *}(c' \cap \beta'_2)) =
h_{1, *}s'_*(c' \cap \beta'_2) =
h_{1, *}(c' \cap (t'_*\beta'_2))
$$
and the final expression computes $c \cap t_*\beta_2$ by construction:
$t'_*\beta'_2 \in \CH_k(Y'_1)$ is a class whose image by $g_{1, *}$ is
$t_*\beta_2$. This proves condition (1).
The other conditions are proved in the same manner
and we omit the detailed arguments.
\end{proof}















